\section{Compiler passes}
\label{sec:appx-compiler}

Section~\ref{sec:design-compiler} summarizes how \sys{} derives
bubble-sized training work from the serving model; this appendix
lists the five passes.
\emph{Operation substitution} replaces each inference-only operation
with a differentiable one over the same weights: paged attention
becomes causal SDPA over the full sequence, MLA expands its latent keys
and values, fused \ac{MoE} becomes grouped \ac{GEMM}, and collectives
get explicit gradient rules. Scaling, positional encoding, routing, and
expert layout are kept, and LoRA deltas are inserted at the adapted
projections.
\emph{Graph capture} traces this forward with
\texttt{AOTAutograd} to obtain the backward graph, with the frozen base
weights as graph inputs so that the fine-tuning worker runs on the live serving
weights.
\emph{Exchange annotation} marks the partial gradients each GPU
produces for cross-GPU summation, with the ranks that take part.
\emph{Partitioning} separates forward from backward and compute from
communication, then uses a roofline cost model to cut compute into
chunks of at most 50\,$\mu$s on the target GPU. Cuts need not follow
layer or operator boundaries: the \ac{MoE} lowering splits the
token--expert pairs of one grouped \ac{GEMM} into several smaller ones.
\emph{Lowering} compiles compute chunks with \texttt{Inductor}, captures
sequences of gated chunks in CUDA graphs to avoid per-chunk launch
cost, and hands communication to \sys{}'s communicator.

\emph{Uneven progress is safe.}
Cross-GPU exchanges use one-sided mailbox messages tagged by exchange and round: a producer publishes without waiting, while consumers wait for the corresponding messages from their peers.
Because every GPU encounters exchanges in the same order and publishing never blocks, uneven progress cannot create a cyclic wait; a faster GPU can only wait until a slower peer reaches the same dependency.
Fine-tuning chains (\S\ref{sec:design-chunk}) synchronize when their gradients are combined for the optimizer step, which bounds how far GPUs drift apart and makes mailbox reuse safe.
Uneven progress may still leave some bubbles unused when a required peer result is not yet ready, but this affects utilization rather than correctness (Appendix~\ref{sec:appx-bubble-use}).

New models built from supported operations need no new code; the same
passes support Qwen1.5-MoE, Qwen3-MoE, DeepSeek-V2-Lite, and OLMoE.
A new operation needs a local lowering and a gradient rule, which can be
tested on its own (Appendix~\ref{appendix:correctness}).
